\centerline{\Huge \bf  Конструктивы в ТЧ}

\begin{flushright}
    \scriptsize{В каждой задаче нужно будет построить пример}
\end{flushright}

\problem{1}
Петя загадал натуральное число $d$ и выписал в ряд $100$ различных натуральных чисел таких, что все разности между соседними числами равны $d$. Могут ли любые два из выписанных чисел быть взаимно просты?

\problem{2}
(а) Существуют ли $100$ натуральных чисел таких, что их сумма равна их наименьшему общему кратному?

\noindent(б) А если числа обязательно различны?

\problem{3}
Существует ли такой набор из $100$ натуральных чисел, что каждое не делится ни на одно из остальных, а квадрат каждого делится на каждое из остальных?

\problem{4}
Существуют ли такие натуральные числа $a_1 < a_2 < a_3 < \ldots < a_{100}$, что 
\[\text{НОД}(a_1, a_2) > \text{НОД}(a_2, a_3) > \ldots > \text{НОД}(a_{99}, a_{100})?\]

\problem{5}
Обозначим через $S(n)$ сумму цифр числа $n$. Найдутся ли три таких натуральных числа $a, b$ и $c$, что 
\[S(a+b) < 5, S(a+c) < 5, S(b+c) < 5,\]
но $S(a+b+c) > 50?$